% !TEX program = xelatex
\documentclass[11pt,a4paper]{ctexart}

% ==================== 页面布局 ====================
\usepackage[top=2.5cm, bottom=2.5cm, left=2.8cm, right=2.8cm]{geometry}
\usepackage{setspace}
\onehalfspacing

% ==================== 字体 ====================
\setCJKmainfont{SimSun}[BoldFont=SimHei]
\setCJKsansfont{SimHei}
\setCJKmonofont{KaiTi}

% ==================== 数学与公式 ====================
\usepackage{amsmath, amssymb, amsthm}
\usepackage{bm}

% ==================== 图表 ====================
\usepackage{graphicx}
\usepackage{float}
\usepackage{booktabs}
\usepackage{longtable}
\usepackage{array}
\usepackage{caption}
\usepackage{subcaption}
\usepackage{xcolor}
\usepackage{colortbl}

% ==================== 代码 ====================
\usepackage{listings}
\lstset{
    basicstyle=\ttfamily\small,
    keywordstyle=\color{blue},
    commentstyle=\color{gray},
    stringstyle=\color{orange},
    numbers=left,
    numberstyle=\tiny\color{gray},
    frame=single,
    breaklines=true,
    breakatwhitespace=false,
    showstringspaces=false,
    tabsize=4,
    language=Python,
    captionpos=b
}

% ==================== 超链接 ====================
\usepackage[colorlinks=true, linkcolor=blue, citecolor=blue, urlcolor=blue]{hyperref}

% ==================== 页眉页脚 ====================
\usepackage{fancyhdr}
\pagestyle{fancy}
\fancyhf{}
\fancyhead[L]{\small CQL 离线强化学习风力发电机桨距角控制}
\fancyhead[R]{\small 技术说明报告}
\fancyfoot[C]{\thepage}
\renewcommand{\headrulewidth}{0.4pt}

% ==================== 标题格式 ====================
\usepackage{titlesec}
\titleformat{\section}{\Large\bfseries\sffamily}{\thesection}{1em}{}
\titleformat{\subsection}{\large\bfseries\sffamily}{\thesubsection}{1em}{}
\titleformat{\subsubsection}{\normalsize\bfseries\sffamily}{\thesubsubsection}{1em}{}

% ==================== 自定义命令 ====================
\newcommand{\pkg}[1]{\texttt{#1}}
\newcommand{\file}[1]{\texttt{#1}}
\newcommand{\code}[1]{\texttt{#1}}
\newcommand{\param}[1]{\textsf{#1}}

% ==================== 文档信息 ====================
\title{
    \vspace{2cm}
    {\huge \textbf{CQL 离线强化学习}}\\[0.3cm]
    {\huge \textbf{风力发电机桨距角控制}}\\[0.8cm]
    {\large \textit{——技术说明报告——}}
}
\author{offline\_pitch\_rl 项目组}
\date{2026年6月20日}

\begin{document}

% ==================== 封面 ====================
\maketitle
\thispagestyle{empty}
\newpage

% ==================== 目录 ====================
\setcounter{page}{1}
\tableofcontents
\newpage

% ==================== 第1章 项目概述 ====================
\section{项目概述}

\subsection{问题背景}

传统风力发电机桨距角控制依赖 PI 控制器。本项目的目标是\textbf{利用离线历史数据训练一个强化学习策略网络（Actor）}，替代 PI 控制器输出桨距角指令。核心挑战在于：

\begin{enumerate}
    \item \textbf{离线数据约束}：无法在线交互，仅能使用 PI 控制器产生的历史数据 (\file{PI\_OfflineData.csv})
    \item \textbf{分布外（OOD）动作过估计}：标准 Q-learning 会对数据分布外的动作产生过高的 Q 值估计
    \item \textbf{策略安全约束}：Actor 输出必须在物理可行范围 $[0, 1.57]$~rad 内，且不能偏离数据分布太远
\end{enumerate}

\subsection{解决方案}

采用 \textbf{CQL (Conservative Q-Learning)} 算法，核心思路：

\begin{enumerate}
    \item \textbf{保守Q值估计}：对 Critic 施加 log-sum-exp 正则化，压低数据分布外随机动作的 Q 值
    \item \textbf{行为克隆正则化}：Actor 损失中加入 BC 项 $\text{MSE}(\pi(s), a_{\text{dataset}})$，防止策略偏离数据分布
    \item \textbf{TD3+BC 自适应权重}：$\lambda = \alpha / \text{mean}(|Q|)$ 动态平衡 Q 最大化与行为克隆
    \item \textbf{Z-score 归一化}：所有状态特征标准化，提升训练数值稳定性
    \item \textbf{奖励缩放}：将原始 $\sim[-200, 0]$ 奖励映射到 $\sim[-2, 0]$
\end{enumerate}

\newpage

% ==================== 第2章 系统架构 ====================
\section{系统架构}

\subsection{整体架构图}

项目采用分层架构设计，自上而下分为配置层、数据层、模型层、训练引擎层、输出部署层和可视化层。

\begin{figure}[H]
\centering
\begin{verbatim}
┌─────────────────────────────────────────────────────────────────┐
│                     offline_pitch_rl 系统架构                      │
├─────────────────────────────────────────────────────────────────┤
│                                                                   │
│  ┌──────────┐    ┌──────────────┐    ┌──────────────────────┐   │
│  │ config.py │◄───│  配置文件     │    │  全局超参数 & 环境定义  │   │
│  └──────────┘    └──────────────┘    └──────────────────────┘   │
│                                                                   │
│  ┌──────────────────────────────────────────────────────────┐    │
│  │                    data/ 数据层                            │    │
│  │  ┌──────────────┐  ┌──────────────────────────────────┐  │    │
│  │  │ dataset.py   │  │ normalizer.py                    │  │    │
│  │  │ · CSV加载     │  │ · Z-score fit/transform/inverse  │  │    │
│  │  │ · 状态构造     │  │ · save/load 持久化               │  │    │
│  │  │ · 奖励计算     │  │                                  │  │    │
│  │  │ · 五元组生成   │  │                                  │  │    │
│  │  └──────────────┘  └──────────────────────────────────┘  │    │
│  └──────────────────────────────────────────────────────────┘    │
│                              │                                    │
│                              ▼                                    │
│  ┌──────────────────────────────────────────────────────────┐    │
│  │                   models/ 模型层                          │    │
│  │  ┌────────────────┐  ┌────────────────────────────────┐  │    │
│  │  │ actor.py       │  │ cql.py (Critic)                │  │    │
│  │  │ · 5→256→256→   │  │ · (5+1)→256→256→128→1         │  │    │
│  │  │   128→1        │  │ · Q(s, a) 值估计               │  │    │
│  │  │ · sigmoid 约束  │  │                                │  │    │
│  │  └────────────────┘  └────────────────────────────────┘  │    │
│  └──────────────────────────────────────────────────────────┘    │
│                              │                                    │
│                              ▼                                    │
│  ┌──────────────────────────────────────────────────────────┐    │
│  │                    train.py 训练引擎                      │    │
│  │  · CQL Critic更新（Bellman + logsumexp保守惩罚）           │    │
│  │  · Actor更新（max Q + BC正则化 + 自适应λ）                 │    │
│  │  · 目标网络软更新（Polyak averaging）                      │    │
│  │  · 梯度裁剪 & 训练历史记录                                 │    │
│  └──────────────────────────────────────────────────────────┘    │
│                              │                                    │
│                              ▼                                    │
│  ┌──────────────────────────────────────────────────────────┐    │
│  │                    输出 & 部署层                           │    │
│  │  ┌───────────────┐  ┌───────────────┐  ┌──────────────┐ │    │
│  │  │ actor_best.pth│  │ actor.onnx    │  │ training_    │ │    │
│  │  │ (PyTorch模型)  │  │ (ONNX部署)    │  │ history.json │ │    │
│  │  └───────────────┘  └───────────────┘  └──────────────┘ │    │
│  └──────────────────────────────────────────────────────────┘    │
│                                                                   │
│  ┌──────────────────────────────────────────────────────────┐    │
│  │                 visualization/ 可视化层                    │    │
│  │  plot_training.py (训练曲线)  │  evaluate.py (评估分析)    │    │
│  └──────────────────────────────────────────────────────────┘    │
│                                                                   │
└─────────────────────────────────────────────────────────────────┘
\end{verbatim}
\caption{系统整体架构图}
\end{figure}

\subsection{数据流}

\begin{figure}[H]
\centering
\begin{verbatim}
PI_OfflineData.csv
       │
       ▼
  dataset.py ─────────────────────────────
  │ 读取CSV                               │
  │ 提取5维状态: [GenSpeed, SpeedError,   │
  │   WindSpeed, PitchAngle, TowerAcc]     │
  │ 构造奖励: -10·err² -2·tower_acc²     │
  │ 缩放奖励: ×0.01                       │
  │ 构造next_state: shift(-1)             │
  │ 构造done: 最后一行为1                  │
  │                                         │
  │ normalizer.fit(state)                  │
  │ state = normalizer.transform(state)    │
  │ next_state = normalizer.transform(ns)  │
  │                                         │
  │ → TensorDataset(s, a, r, ns, d)        │
  └────────────────────────────────────────┘
       │
       ▼
  DataLoader (batch_size=512, shuffle=True)
       │
       ▼
  train.py ── 500 epochs ──► output/
\end{verbatim}
\caption{数据流水线}
\end{figure}

\newpage

% ==================== 第3章 MDP建模 ====================
\section{马尔可夫决策过程建模}

\subsection{状态空间}

状态向量 $\mathbf{s} \in \mathbb{R}^5$，五个维度如下：

\begin{table}[H]
\centering
\caption{状态空间定义}
\begin{tabular}{cclll}
\toprule
\textbf{维度} & \textbf{变量名} & \textbf{含义} & \textbf{单位} & \textbf{归一化} \\
\midrule
0 & GenSpeed    & 发电机转速       & -- & Z-score \\
1 & SpeedError  & 转速误差 (Ref $-$ GenSpeed) & -- & Z-score \\
2 & WindSpeed   & 风速             & -- & Z-score \\
3 & PitchAngle  & 当前桨距角测量值  & -- & Z-score \\
4 & TowerAcc    & 塔架加速度       & -- & Z-score \\
\bottomrule
\end{tabular}
\end{table}

状态维度 $\texttt{STATE\_DIM} = 5$，所有状态经 Z-score 归一化（均值0，标准差1）。

\subsection{动作空间}

\begin{table}[H]
\centering
\caption{动作空间定义}
\begin{tabular}{ll}
\toprule
\textbf{属性} & \textbf{值} \\
\midrule
维度   & 1 ($\texttt{ACTION\_DIM} = 1$) \\
含义   & 桨距角指令 (Pitch Demand) \\
范围   & $[\texttt{MIN\_PITCH}, \texttt{MAX\_PITCH}] = [0.0, 1.57]$~rad \\
\bottomrule
\end{tabular}
\end{table}

Actor 输出通过 $\sigma(\cdot)$ (sigmoid) 映射到 $[0, 1]$，再缩放到 $[0, 1.57]$：
\begin{equation}
    a = 0.0 + \sigma(f_{\theta}(\mathbf{s})) \times (1.57 - 0.0)
\end{equation}

\subsection{奖励函数}

即时奖励由转速误差和塔架加速度两项惩罚组成：

\begin{equation}
    r_{\text{raw}} = -10.0 \cdot (\text{speed\_error})^2 - 2.0 \cdot (\text{tower\_acc})^2
\end{equation}
\begin{equation}
    r = r_{\text{raw}} \times 0.01 \quad (\text{缩放至 } \sim[-2, 0])
\end{equation}

\textbf{设计意图}：
\begin{itemize}
    \item 转速误差惩罚（权重 10）：约束发电机转速跟随参考值
    \item 塔架加速度惩罚（权重 2）：抑制塔架振动，保护机械结构
    \item 奖励缩放（因子 0.01）：将原始范围 $\sim[-200, 0]$ 映射到 $\sim[-2, 0]$，避免大梯度破坏训练稳定性
\end{itemize}

\subsection{转移与终止}

\begin{itemize}
    \item \textbf{下一状态} $\mathbf{s}_{t+1}$：由于是离线数据，采用 $\text{shift}(-1)$ 近似——即 $\mathbf{s}[t+1]$ 作为 $\mathbf{s}[t]$ 的下一状态，最后一行的 $\mathbf{s}_{t+1}$ 设为自身
    \item \textbf{终止标志} $d$：全零向量，仅最后一行设为 1，表示轨迹结束
    \item \textbf{折扣因子}：$\gamma = 0.99$
\end{itemize}

\newpage

% ==================== 第4章 数据流水线 ====================
\section{数据流水线}

\subsection{数据加载}

核心函数 \code{load\_dataset(csv\_file)} 位于 \file{data/dataset.py}，执行以下流程：

\begin{enumerate}
    \item 使用 \pkg{pandas} 读取 CSV 文件
    \item 提取五列原始数据：GenSpeed, GenSpeed\_Ref, WindSpeed, PitchAngle\_Mea, TowerAcc
    \item 构造转速误差：$\text{SpeedError} = \text{GenSpeed\_Ref} - \text{GenSpeed}$
    \item 构造动作：直接使用 Action\_PI\_Dem 列
    \item 计算奖励：$r = (-10.0 \cdot e_{\text{speed}}^2 - 2.0 \cdot a_{\text{tower}}^2) \times 0.01$
    \item 构造下一状态：$\text{next\_state}[i] = \text{state}[i+1]$, $\text{next\_state}[-1] = \text{state}[-1]$
    \item Z-score 归一化：\code{normalizer.fit(state)} → \code{normalizer.transform(state)}, \code{normalizer.transform(next\_state)}
    \item 返回五元组 \code{(s, a, r, ns, d, norm)}
\end{enumerate}

\subsection{Z-score 归一化器}

\file{data/normalizer.py} 中的 \code{Normalizer} 类实现 Z-score 标准化：

\begin{table}[H]
\centering
\caption{Normalizer 类方法}
\begin{tabular}{ll}
\toprule
\textbf{方法} & \textbf{功能} \\
\midrule
\code{fit(x)}     & 计算各维度均值 $\mu$ 和标准差 $\sigma$；若 $\sigma < 10^{-8}$ 则置为 1.0（防零除） \\
\code{transform(x)} & $(x - \mu) / \sigma$ \\
\code{inverse(x)}   & $x \cdot \sigma + \mu$ \\
\code{save(path)}   & 保存 $\mu$/$\sigma$ 到 JSON 文件 \\
\code{load(path)}   & 从 JSON 加载 $\mu$/$\sigma$（类方法） \\
\bottomrule
\end{tabular}
\end{table}

归一化参数在训练时保存至 \file{output/normalization.json}，供推理部署时使用。

\newpage

% ==================== 第5章 神经网络模型 ====================
\section{神经网络模型}

\subsection{Actor 策略网络}

\file{models/actor.py} 定义策略网络 $\pi_{\theta}(\mathbf{s})$，结构如下：

\begin{verbatim}
输入: state [B, 5]
  │
  ▼ Linear(5, 256) + ReLU
  │
  ▼ Linear(256, 256) + ReLU
  │
  ▼ Linear(256, 128) + ReLU
  │
  ▼ Linear(128, 1)
  │
  ▼ sigmoid → [0, 1]
  │
  ▼ scale: MIN_PITCH + x * (MAX_PITCH - MIN_PITCH)
  │
输出: pitch_demand [B, 1] ∈ [0, 1.57]
\end{verbatim}

\textbf{设计要点}：
\begin{itemize}
    \item 三层隐藏层（256→256→128），使用 ReLU 激活函数
    \item 输出层不做激活，而是通过 sigmoid 缩放至合法动作范围
    \item 参数量约 101,633
\end{itemize}

\subsection{Critic Q值网络}

\file{models/cql.py} 定义 Q 值网络 $Q_{\phi}(\mathbf{s}, a)$：

\begin{verbatim}
输入: state [B, 5] + action [B, 1] = [B, 6]
  │
  ▼ concat(s, a) along dim=-1
  │
  ▼ Linear(6, 256) + ReLU
  │
  ▼ Linear(256, 256) + ReLU
  │
  ▼ Linear(256, 128) + ReLU
  │
  ▼ Linear(128, 1)
  │
输出: Q(s, a) [B, 1]
\end{verbatim}

\textbf{设计要点}：
\begin{itemize}
    \item 状态与动作在最后一维拼接后送入网络
    \item 与 Actor 结构对称（隐藏层维度一致），便于训练平衡
    \item 目标网络 $\bar{Q}_{\phi'}$ 为 Critic 的 Polyak 平均副本
\end{itemize}

\newpage

% ==================== 第6章 CQL训练算法 ====================
\section{CQL 训练算法}

\subsection{训练循环总览}

训练主循环位于 \file{train.py}，共 500 个 epoch，每个 batch 执行三步更新：

\begin{enumerate}
    \item \textbf{Critic 更新}：最小化 Bellman 误差 + CQL 保守性惩罚
    \item \textbf{Actor 更新}：最大化 Q 值 + 行为克隆正则化（TD3+BC 自适应 $\lambda$）
    \item \textbf{目标网络软更新}：Polyak 平均
\end{enumerate}

\subsection{Critic 更新}

\subsubsection{Bellman 误差（标准 TD 学习）}

\begin{align}
    \tilde{a} &= \pi(\mathbf{s}_{t+1}) \tag{目标策略动作} \\
    \tilde{q} &= \bar{Q}(\mathbf{s}_{t+1}, \tilde{a}) \tag{目标Q值} \\
    y &= r + \gamma \cdot (1 - d) \cdot \tilde{q} \tag{TD目标} \\
    \mathcal{L}_{\text{bellman}} &= \frac{1}{B}\sum_{i=1}^{B} \left(Q(\mathbf{s}_i, a_i) - y_i\right)^2 \tag{Bellman MSE}
\end{align}

\subsubsection{CQL 保守性惩罚}

对每批次数据采样 $N=10$ 个均匀随机动作，在动作维度上计算 log-sum-exp：

\begin{align}
    a_{\text{rand}} &\sim \mathcal{U}(0, \texttt{MAX\_PITCH})^{B \times N \times 1} \tag{随机动作} \\
    Q_{\text{rand}} &= Q(\mathbf{s}, a_{\text{rand}}) \quad \in \mathbb{R}^{B \times N} \tag{随机动作Q值} \\
    \mathcal{L}_{\text{cql}} &= \frac{1}{B}\sum_{i=1}^{B} \log\sum_{j=1}^{N} \exp Q_{\text{rand}}[i, j] \;-\; \frac{1}{B}\sum_{i=1}^{B} Q(\mathbf{s}_i, a_i) \tag{CQL惩罚}
\end{align}

\textbf{直觉解释}：
\begin{itemize}
    \item log-sum-exp 项推高所有随机动作的 Q 值 → 优化时该项降低 → 随机动作 Q 值被压低
    \item $-\text{mean}(Q)$ 项确保数据集中真实动作的 Q 值不被压低
\end{itemize}

\subsubsection{Critic 总损失}

\begin{equation}
    \mathcal{L}_{\text{critic}} = \mathcal{L}_{\text{bellman}} + \alpha_{\text{cql}} \cdot \mathcal{L}_{\text{cql}}
\end{equation}

其中 $\alpha_{\text{cql}} = 1.0$ 控制保守性强度。

\subsection{Actor 更新：TD3+BC 自适应行为克隆}

\begin{align}
    q_{\text{pred}} &= Q(\mathbf{s}, \pi(\mathbf{s})) \tag{当前策略Q值} \\
    \mathcal{L}_{\text{bc}} &= \frac{1}{B}\sum_{i=1}^{B} (\pi(\mathbf{s}_i) - a_i)^2 \tag{行为克隆损失} \\
    \lambda &= \frac{\lambda_{\text{bc}}}{\text{mean}(|Q|) + 10^{-6}} \tag{自适应权重} \\
    \mathcal{L}_{\text{actor}} &= -\lambda \cdot \frac{1}{B}\sum_{i=1}^{B} q_{\text{pred},i} \;+\; \mathcal{L}_{\text{bc}} \tag{Actor总损失}
\end{align}

\textbf{自适应 $\lambda$ 机制}：当 Q 值幅度大时自动降低 BC 权重，当 Q 值幅度小时增大 BC 权重——保证 BC 与 Q 最大化的稳健平衡。

\subsection{目标网络软更新}

\begin{equation}
    \phi' \leftarrow \tau \cdot \phi + (1 - \tau) \cdot \phi', \quad \tau = 0.005
\end{equation}

Polyak 平均使目标网络平滑变化，稳定 TD 学习目标。

\subsection{训练稳定性措施}

\begin{table}[H]
\centering
\caption{训练稳定性措施汇总}
\begin{tabular}{lll}
\toprule
\textbf{措施} & \textbf{参数/方式} & \textbf{目的} \\
\midrule
梯度裁剪   & \texttt{GRAD\_CLIP = 10.0}   & 防止梯度爆炸 \\
奖励缩放   & \texttt{REWARD\_SCALE = 0.01} & 数值稳定 \\
状态归一化 & Z-score                  & 消除量纲差异 \\
降低保守性 & \texttt{CQL\_ALPHA = 1.0}（原 5.0） & 避免Q值过度悲观 \\
学习率一致 & \texttt{LR\_CRITIC = 1e-4}   & 训练平衡 \\
\bottomrule
\end{tabular}
\end{table}

\newpage

% ==================== 第7章 超参数与设计决策 ====================
\section{关键设计决策与超参数}

\subsection{超参数总表}

\begin{table}[H]
\centering
\caption{全局超参数配置}
\begin{tabular}{lll}
\toprule
\textbf{参数} & \textbf{值} & \textbf{说明} \\
\midrule
\texttt{STATE\_DIM}      & 5      & 状态维度 \\
\texttt{ACTION\_DIM}     & 1      & 动作维度 \\
\texttt{MIN\_PITCH}      & 0.0    & 最小桨距角 (rad) \\
\texttt{MAX\_PITCH}      & 1.57   & 最大桨距角 (rad) \\
\texttt{BATCH\_SIZE}     & 512    & 批大小 \\
\texttt{GAMMA}           & 0.99   & 折扣因子 \\
\texttt{TAU}             & 0.005  & 目标网络软更新系数 \\
\texttt{LR\_ACTOR}       & $10^{-4}$ & Actor 学习率 \\
\texttt{LR\_CRITIC}      & $10^{-4}$ & Critic 学习率 \\
\texttt{CQL\_ALPHA}      & 1.0    & CQL 保守性惩罚权重 \\
\texttt{CQL\_NUM\_SAMPLES}& 10     & 每状态随机动作采样数 \\
\texttt{LAMBDA\_BC}      & 2.5    & BC 正则化基础权重 \\
\texttt{REWARD\_SCALE}   & 0.01   & 奖励缩放因子 \\
\texttt{GRAD\_CLIP}      & 10.0   & 梯度裁剪阈值 \\
\texttt{EPOCHS}          & 500    & 训练轮数 \\
\texttt{DEVICE}          & cuda   & 计算设备 \\
\bottomrule
\end{tabular}
\end{table}

\subsection{关键设计取舍}

\begin{table}[H]
\centering
\caption{设计决策理由分析}
\begin{tabular}{p{3.5cm}p{8.5cm}}
\toprule
\textbf{设计选择} & \textbf{理由} \\
\midrule
CQL over BCQ/BEAR & CQL 实现简单，log-sum-exp 惩罚在动作维度上直接且有效 \\

TD3+BC 自适应 $\lambda$ over 固定 $\lambda$ & Q 值幅度随训练变化，固定权重难以平衡 Q 最大化与 BC 正则化 \\

Z-score over MinMax 归一化 & Z-score 对异常值更鲁棒，保持分布形状 \\

单 Critic over 双 Critic & 简化实现；CQL 本身的保守性已提供足够的低估偏置 \\

ONNX opset 11 & 广泛兼容的部署标准 \\

批大小 512 over 256 & 离线数据充足时，大批量提供更稳定的梯度估计 \\
\bottomrule
\end{tabular}
\end{table}

\newpage

% ==================== 第8章 模型导出与部署 ====================
\section{模型导出与部署}

\subsection{ONNX 导出}

\file{export\_onnx.py} 将训练好的 Actor 模型导出为 ONNX 格式：

\begin{lstlisting}[caption={ONNX 导出流程}, label={lst:onnx}]
model = Actor(STATE_DIM)
model.load_state_dict(torch.load("output/actor_best.pth", map_location="cpu"))
model.eval()

dummy = torch.randn(1, STATE_DIM)

torch.onnx.export(
    model, dummy, "output/actor.onnx",
    input_names=["state"],
    output_names=["pitch_demand"],
    dynamic_axes={
        "state": {0: "batch"},
        "pitch_demand": {0: "batch"}
    },
    opset_version=11
)
\end{lstlisting}

\textbf{导出配置}：
\begin{itemize}
    \item \texttt{input\_names}: \code{["state"]}
    \item \texttt{output\_names}: \code{["pitch\_demand"]}
    \item \texttt{dynamic\_axes}: batch 维度设为动态，支持任意批次大小
    \item \texttt{opset\_version}: 11
\end{itemize}

\subsection{推理部署所需文件}

\begin{table}[H]
\centering
\caption{部署文件清单}
\begin{tabular}{ll}
\toprule
\textbf{文件} & \textbf{用途} \\
\midrule
\file{output/actor.onnx}        & ONNX 模型（跨平台推理） \\
\file{output/actor\_best.pth}   & PyTorch 权重（可继续训练） \\
\file{output/normalization.json}& 归一化参数（推理时需对输入做相同归一化） \\
\bottomrule
\end{tabular}
\end{table}

\newpage

% ==================== 第9章 可视化与评估 ====================
\section{可视化与评估体系}

\subsection{训练过程监控}

\file{visualization/plot\_training.py} 生成三类图表，所有曲线包含原始值和 10-epoch 滑动平均平滑线，支持中英双语标注。

\begin{table}[H]
\centering
\caption{训练监控图表}
\begin{tabular}{lll}
\toprule
\textbf{图表} & \textbf{内容} & \textbf{输出文件} \\
\midrule
Loss Curves      & Actor/Critic/Bellman/CQL 四合一损失曲线 & \file{training\_losses.png} \\
Q-Value \& Return & 平均 Q 值趋势 \& 平均即时奖励          & \file{q\_value\_return.png} \\
Training Summary & 6合1综合训练摘要图                       & \file{training\_summary.png} \\
\bottomrule
\end{tabular}
\end{table}

\subsection{策略评估}

\file{visualization/evaluate.py} 提供四维度策略质量评估：

\begin{table}[H]
\centering
\caption{评估分析维度}
\begin{tabular}{lll}
\toprule
\textbf{评估维度} & \textbf{分析方法} & \textbf{输出文件} \\
\midrule
动作对比   & 散点图(pred vs true)、误差直方图、时序对比、分布对比 & \file{action\_comparison.png} \\
策略质量   & 动作偏差分布、逐样本偏差、状态-动作关系           & \file{policy\_quality.png} \\
状态覆盖   & 5维状态 vs 动作散点图、动作值分布                 & \file{state\_action\_coverage.png} \\
\bottomrule
\end{tabular}
\end{table}

\subsubsection{定量指标}

评估输出以下指标：
\begin{itemize}
    \item MAE (Mean Absolute Error)：$\frac{1}{N}\sum |\pi(\mathbf{s}_i) - a_i|$
    \item RMSE (Root Mean Square Error)：$\sqrt{\frac{1}{N}\sum (\pi(\mathbf{s}_i) - a_i)^2}$
    \item Max Absolute Error：$\max_i |\pi(\mathbf{s}_i) - a_i|$
    \item 动作均值/标准差对比
    \item 奖励统计
\end{itemize}

\subsubsection{评估结论等级}

\begin{table}[H]
\centering
\caption{MAE 等级判定}
\begin{tabular}{cl}
\toprule
\textbf{MAE 范围} & \textbf{等级} \\
\midrule
$< 0.01$ & 优秀 (EXCELLENT) — Actor 紧密匹配行为策略 \\
$< 0.05$ & 良好 (GOOD) — 与行为策略有轻微偏差 \\
$< 0.10$ & 一般 (FAIR) — 明显偏差，需检查训练 \\
$\geq 0.10$ & 较差 (POOR) — 偏差较大，训练可能需要调参 \\
\bottomrule
\end{tabular}
\end{table}

\newpage

% ==================== 第10章 模块清单 ====================
\section{模块清单}

\begin{table}[H]
\centering
\caption{项目模块一览}
\begin{tabular}{p{5.5cm}p{6.5cm}r}
\toprule
\textbf{文件路径} & \textbf{职责} & \textbf{行数} \\
\midrule
\file{config.py}                  & 全局超参数配置                          & 29 \\
\file{train.py}                   & CQL 训练主循环                          & 241 \\
\file{export\_onnx.py}            & ONNX 模型导出                           & 47 \\
\file{data/dataset.py}           & CSV 加载、状态/奖励构造、五元组生成     & 83 \\
\file{data/normalizer.py}        & Z-score 归一化器（含持久化）            & 46 \\
\file{models/actor.py}           & Actor 策略网络 $\pi_{\theta}$           & 40 \\
\file{models/cql.py}             & Critic Q值网络 $Q_{\phi}$               & 33 \\
\file{env/reward.py}             & 奖励函数占位（实际在 \file{dataset.py}） & 33 \\
\file{utils/save\_scaler.py}     & 归一化器保存占位（已集成至 Normalizer） & 2 \\
\file{utils/load\_scaler.py}     & 归一化器加载占位（已集成至 Normalizer） & 2 \\
\file{visualization/plot\_training.py} & 训练曲线可视化（中英双语）         & 208 \\
\file{visualization/evaluate.py}      & 策略评估与可视化（中英双语）       & 268 \\
\bottomrule
\end{tabular}
\end{table}

\newpage

% ==================== 附录 ====================
\appendix
\section{运行命令}

\begin{lstlisting}[caption={运行命令}, label={lst:run}]
# 1. 训练
python train.py

# 2. 生成训练曲线
python visualization/plot_training.py

# 3. 策略评估
python visualization/evaluate.py

# 4. ONNX 导出
python export_onnx.py
\end{lstlisting}

\section{依赖环境}

项目所需 Python 包：
\begin{verbatim}
    torch, numpy, pandas, matplotlib, onnx
\end{verbatim}

\section{CQL 损失函数公式总结}

\subsection*{Critic 总损失}

\begin{equation}
\boxed{
\mathcal{L}_{\text{critic}} = \underbrace{\mathbb{E}_{(s,a,s')\sim\mathcal{D}}\left[\left(Q(s,a) - \left(r + \gamma (1-d) \bar{Q}(s', \pi(s'))\right)\right)^2\right]}_{\text{Bellman TD Error}}
\;+\;
\alpha \cdot \underbrace{\left[\mathbb{E}_{s\sim\mathcal{D}}\left[\log\sum_{a'}\exp Q(s,a')\right] - \mathbb{E}_{(s,a)\sim\mathcal{D}}\left[Q(s,a)\right]\right]}_{\text{CQL Conservative Penalty}}
}
\end{equation}

\subsection*{Actor 总损失}

\begin{equation}
\boxed{
\mathcal{L}_{\text{actor}} = -\underbrace{\frac{\lambda_{\text{bc}}}{\text{mean}(|Q|)}}_{\text{自适应}\lambda} \cdot \mathbb{E}_{s\sim\mathcal{D}}\left[Q(s, \pi(s))\right]
\;+\;
\underbrace{\mathbb{E}_{(s,a)\sim\mathcal{D}}\left[(\pi(s) - a)^2\right]}_{\text{BC Regularization}}
}
\end{equation}

\subsection*{目标网络更新}

\begin{equation}
\boxed{
\phi' \leftarrow \tau \cdot \phi + (1 - \tau) \cdot \phi'
}
\end{equation}

\vspace{2cm}

\begin{center}
\rule{0.5\textwidth}{0.5pt}\\[0.5cm]
{\large \textbf{—— 文档结束 ——}}
\end{center}

\end{document}
